\chapter{Implementation}
\label{ch:impl}

\section{Software structure}

The system is a Python package, \code{qrcav}, with a strict layering: lower
layers never import upper ones, and the protocol code performs no I/O.

\begin{table}[h]
\centering
\caption{Package structure}
\label{tab:modules}
\begin{tabular}{@{}p{4.9cm}p{9.1cm}@{}}
\toprule
Module & Responsibility \\
\midrule
\code{crypto/suites.py} & ML-KEM and ML-DSA behind one interface (OpenSSL); ECDH/ECDSA baseline \\
\code{crypto/kdf.py} & SHAKE-256 hash, HKDF-SHA3-256, session key derivation \\
\code{crypto/records.py} & AES-256-GCM records, counters, replay window \\
\code{codec.py} & canonical compact binary encoding (signatures need identical bytes on both sides) \\
\code{pki.py} & Trusted Authority, certificates, CRL, registry database, identities \\
\code{handshake.py} & the four-flight handshake as sans-IO state machines \\
\code{certcache.py} & certificate references and the bounded LRU cache \\
\code{consensus/agspbft.py} & AGS-PBFT and PBFT replica: normal case, scoring, regrouping, view change \\
\code{consensus/sim.py} & discrete-event harness running real replicas with real signatures \\
\code{ledger.py} & hash-linked, certified blocks; full-chain verification \\
\code{net/secure.py} & secure channels over TCP with length framing \\
\code{nodes/} & one process per entity: TA, CAV, RSU, MEC, consensus node, cloud \\
\code{sim/} & ns-3 scenario (C++), its Python driver, measurement profiles \\
\code{eval/} & benchmarks, attack harness, consensus experiments, ns-3 sweeps, figures, replay viewer \\
\bottomrule
\end{tabular}
\end{table}

\section{Cryptographic suites}

A \emph{suite} pairs a KEM with a signature scheme. Five are implemented:
\textbf{DECK} (ML-KEM-1024 + ML-DSA-65, the pairing in our Review~1 design and
in~\cite{olushola}), \textbf{L1} (512/44), \textbf{L3} (768/65), \textbf{L5}
(1024/87), and \textbf{CLASSIC} (ECDH P-256 + ECDSA P-256, as in IEEE~1609.2).
DECK pairs a category-5 KEM with a category-3 signature, so its overall
strength is category~3, the same as L3.

ECDH is wrapped as a KEM with the DHKEM construction of RFC~9180~\cite{hpke}:
\textsf{Encaps} generates an ephemeral key pair, computes the Diffie--Hellman
value with the receiver's public key, hashes it together with both public
keys, and sends the ephemeral public key as the ``ciphertext''. The classical
baseline therefore runs through exactly the same handshake code as the
post-quantum suites, which makes the comparison like-for-like.

All algorithms come from the \code{cryptography} library (version 47 or
later), which bundles OpenSSL~3.5 with native, constant-time ML-KEM and
ML-DSA, as a deployment would use. One exception: the library exposes
ML-KEM-768 and ML-KEM-1024 but not ML-KEM-512, so the L1 suite takes its KEM
from \code{kyber-py}, a pure-Python implementation of FIPS~203. It produces
the same sizes and encodings, but it is slower and, as its authors state, not
constant-time, which is why L1's handshake is the slowest in
Section~\ref{ch:results}.

\section{Distributed system}

Every entity is a separate operating-system process started from one
configuration file. A run of the demonstration scenario starts 19 processes:
the TA, the cloud, one MEC, eight consensus nodes, three RSUs and five
vehicles. Transport is TCP with a 4-byte length prefix per frame; the four
handshake flights are frames, and afterwards every frame is one AES-GCM
record. Registration and CRL retrieval use plain TCP, which is safe because
certificates and CRLs are TA-signed and verified on receipt. Each process
writes structured JSON logs, from which an end-to-end report is produced.

The configuration injects faults: a consensus node can return wrong results
(\code{wrong\_result}) or stay silent, a node can crash at a given time, the
TA can revoke an entity, and vehicles follow a script of events, V2V messages,
handovers and reconnections.

\section{ns-3 simulation}

\subsection{Method}

ns-3 3.44~\cite{ns3} simulates the radio, MAC, backhaul and mobility. The
scenario is written in C++ and compiled at run time through the ns-3 Python
bindings. ns-3 does not execute the Python cryptography, which would mix
wall-clock and simulated time. Instead the real protocol code is run first to
produce a \emph{measurement profile}: the exact byte length of every message
(the four handshake flights, cached flights, AEAD and signed beacons, events,
alerts) and empirical samples of the processing time at each handshake stage,
per beacon, and for a consensus commit. The simulation draws from these
samples whenever a node processes a message, and each node has a
single-server CPU queue.

\begin{table}[h]
\centering
\caption{ns-3 scenario}
\label{tab:ns3}
\begin{tabular}{@{}p{4.2cm}p{9.8cm}@{}}
\toprule
Parameter & Value \\
\midrule
Radio & IEEE 802.11p, channel 172 (5.86\,GHz), 10\,MHz, 6\,Mb/s, 23\,dBm \\
Propagation & log-distance ($n = 2.2$) + Nakagami fading, $m = 3 / 1.5 / 1$ below 50 / 150 / beyond~\cite{torrent} \\
Detection, CCA & $-85$\,dBm (the 10\,MHz value) \\
Area & 1.5\,km $\times$ 1.5\,km Manhattan grid, 250\,m blocks \\
RSUs & at intersections every 500\,m: 9 RSUs, wired to one MEC (1\,Gb/s, 2\,ms) \\
Vehicles & 50--200; 8--17\,m/s; straight 50\,\%, left or right 25\,\% each at intersections \\
Beacons & 10\,Hz link-layer broadcast (no ACK, no retry) \\
Fragmentation & 1,400-byte chunks; reassembly at the receiver \\
Handshake & 200\,ms timeout, 3 retries; NACK for missing ServerHello chunks \\
Association & start once $\ge 6$ of 10 announcements per second are heard; handover to an RSU heard $\ge 4$ per second better \\
EDCA & beacons, events, alerts as AC\_VO; handshake flights as AC\_BE \\
Backoff & after a failed handshake wait $U(0.5, 1) \times 2^k$\,s \\
Run & 30\,s simulated, first 5\,s excluded; 3 seeds \\
\bottomrule
\end{tabular}
\end{table}

\subsection{Simulator pitfalls found and corrected}

Five ns-3 defaults distorted results by an order of magnitude before they
were corrected, and they belong in any reproduction of this work:
\begin{enumerate}
\item \textbf{ARP.} The ARP cache queues only three packets per unresolved
address, so most fragments of a first ServerHello were silently dropped.
802.11p safety messaging has no ARP, so neighbour caches are pre-populated.
\item \textbf{Preamble detection.} ns-3 defaults to the 20\,MHz value of
$-82$\,dBm; 802.11p's 10\,MHz channel uses $-85$\,dBm. The wrong value cut
coverage from 99\,\% to about 80\,\%.
\item \textbf{Unicast beacons} suffered up to 7 MAC retries each; real beacons
are broadcasts.
\item \textbf{Whole-flight retransmission} after one lost fragment caused a
retransmission storm; it was replaced with selective retransmission.
\item \textbf{First-contact association.} Starting a 17-frame handshake on the
first announcement heard, usually at the edge of coverage, gave 5\,\%
success; association is now gated on link quality.
\end{enumerate}

\subsection{Certificate caching in the simulation}

Mode 1 models the deployment of Section~\ref{sec:cache}: RSU certificates are
pre-distributed, so every ServerHello carries a reference; after its first
session a vehicle sends a reference too, and an RSU that has not seen the
vehicle resolves it from the MEC directory at the cost of one backhaul round
trip (no airtime). Mode 2, used as an upper bound, starts with all caches warm.

\subsection{Replay viewer}

For demonstration, the simulation can record an event trace (handshake
start, session up, failure, link loss, handover, event, alert) and the
channel airtime in 0.5\,s buckets. A script runs the same scenario (same
mobility and seed) for the classical suite, DECK and DECK with caching, and
writes a self-contained web page that replays the three runs on the city grid
(Figure~\ref{fig:viewer}).

\section{Evaluation harness}

\begin{itemize}
\item \textbf{Primitives}: 200 timed runs per operation, after warm-up runs
are discarded.
\item \textbf{Handshakes}: 100 complete handshakes per suite in memory,
recording each stage's processing time and every flight's size, with and
without caching.
\item \textbf{Beacon budget}: per-beacon cryptography plus airtime for AEAD
records and for signed broadcasts.
\item \textbf{Attacks}: \V{attacks-total} attacks run against the real
implementation; an attack counts as rejected only if the protocol aborts or
drops the message.
\item \textbf{Consensus}: real replicas with real ML-DSA signatures in a
discrete-event loop, for 8--24 nodes, AGS-PBFT with and without shadow
validation against PBFT; fault tolerance; view change; score dynamics.
\item \textbf{ns-3 sweeps}: five suites and DECK with caching at 50--200
vehicles over three seeds, plus single-seed variants.
\end{itemize}
Timing measurements were taken with the machine otherwise idle, because
benchmarks run alongside ns-3 inflated latencies.

\section{Testing}

55 automated tests (\code{pytest}) cover the codec, FIPS key and signature sizes, key derivation (output lengths and transcript binding), record
counters and replay window, ML-KEM-512 sizes, registration (proof of
possession, roster, duplicate keys), handshakes for all suites, role checks,
certificate caching (savings, fallback, revocation, substituted references,
forged HelloRetry), AGS-PBFT (normal case, scoring, regrouping, view change,
Byzantine nodes, determinism) and one complete multi-process run of the
distributed system.

\section{Bugs found by our own verification}

\begin{enumerate}
\item \textbf{Consensus stalled about once in 14 runs.} Nodes checked whether
the reporting vehicle was registered against the registry copy each had
fetched at start-up; a node that started before the vehicle enrolled
disagreed with the rest, and with one Byzantine node in a set of four no
result reached $2f+1$. Found by tracing every message of a failing run, fixed
by carrying the reporter's TA-signed certificate in the request so validation
is a pure function of the request, and verified with 0 failures in 25 runs.
\item \textbf{Sequence-number reuse after a view change}: the new leader
re-issued prepared requests and then reused their numbers.
\item \textbf{PRE-PREPAREs overtaking NEW-VIEW} were dropped as wrong-view;
they are now buffered and replayed.
\item \textbf{Measurement contention}: timings taken while ns-3 used both
cores were inflated and were re-measured.
\end{enumerate}
